\section{The theorem on formal functions}
\label{section-theorem-formal-functions}

\noindent
In this section we study the behaviour of cohomology of
sequences of sheaves either of the form $\{I^n\mathcal{F}\}_{n \geq 0}$
or of the form $\{\mathcal{F}/I^n\mathcal{F}\}_{n \geq 0}$ as
$n$ varies.

\medskip\noindent
Here and below we use the following notation.
Given a morphism of schemes $f : X \to Y$, a quasi-coherent sheaf
$\mathcal{F}$ on $X$, and a quasi-coherent sheaf of ideals
$\mathcal{I} \subset \mathcal{O}_Y$ we denote
$\mathcal{I}^n\mathcal{F}$ the quasi-coherent subsheaf generated
by products of local sections of $f^{-1}(\mathcal{I}^n)$ and
$\mathcal{F}$. In a formula
$$
\mathcal{I}^n\mathcal{F}
=
\Im\left(
f^*(\mathcal{I}^n) \otimes_{\mathcal{O}_X} \mathcal{F}
\longrightarrow
\mathcal{F}
\right).
$$
Note that there are natural maps
$$
f^{-1}(\mathcal{I}^n) \otimes_{f^{-1}\mathcal{O}_Y} \mathcal{I}^m\mathcal{F}
\longrightarrow
f^*(\mathcal{I}^n) \otimes_{\mathcal{O}_X} \mathcal{I}^m\mathcal{F}
\longrightarrow
\mathcal{I}^{n + m}\mathcal{F}
$$
Hence a section of $\mathcal{I}^n$ will give rise to a
map $R^pf_*(\mathcal{I}^m\mathcal{F}) \to
R^pf_*(\mathcal{I}^{n + m}\mathcal{F})$ by functoriality
of higher direct images. Localizing and then sheafifying we
see that there are $\mathcal{O}_Y$-module maps
$$
\mathcal{I}^n \otimes_{\mathcal{O}_Y} R^pf_*(\mathcal{I}^m\mathcal{F})
\longrightarrow
R^pf_*(\mathcal{I}^{n + m}\mathcal{F}).
$$
In other words we see that
$\bigoplus_{n \geq 0} R^pf_*(\mathcal{I}^n\mathcal{F})$
is a graded $\bigoplus_{n \geq 0} \mathcal{I}^n$-module.

\medskip\noindent
If $Y = \Spec(A)$ and $\mathcal{I} = \widetilde{I}$ we denote
$\mathcal{I}^n\mathcal{F}$ simply $I^n\mathcal{F}$. The maps
introduced above give $M = \bigoplus H^p(X, I^n\mathcal{F})$ the
structure of a graded $S = \bigoplus I^n$-module. If $f$ is proper,
$A$ is Noetherian and $\mathcal{F}$ is coherent, then this turns out
to be a module of finite type.

\begin{lemma}
\label{lemma-cohomology-powers-ideal-times-F}
\begin{reference}
\cite[III Cor 3.3.2]{EGA}
\end{reference}
Let $A$ be a Noetherian ring.
Let $I \subset A$ be an ideal.
Set $B = \bigoplus_{n \geq 0} I^n$.
Let $f : X \to \Spec(A)$ be a proper morphism.
Let $\mathcal{F}$ be a coherent sheaf on $X$.
Then for every $p \geq 0$ the graded $B$-module
$\bigoplus_{n \geq 0} H^p(X, I^n\mathcal{F})$ is
a finite $B$-module.
\end{lemma}

\begin{proof}
Let $\mathcal{B} = f^*\widetilde{B}$. Since $\mathcal{B}$ surjects
onto $\bigoplus I^n\mathcal{O}_X$ it is clear that
$\bigoplus I^n\mathcal{F}$ is a finite type graded $\mathcal{B}$-module.
Hence the result follows from Lemma \ref{lemma-graded-finiteness} part (1).
\end{proof}

\begin{lemma}
\label{lemma-cohomology-powers-ideal-times-sheaf}
Given a morphism of schemes $f : X \to Y$, a quasi-coherent sheaf
$\mathcal{F}$ on $X$, and a quasi-coherent sheaf of ideals
$\mathcal{I} \subset \mathcal{O}_Y$. Assume $Y$ locally
Noetherian, $f$ proper, and $\mathcal{F}$ coherent.
Then
$$
\mathcal{M} =
\bigoplus\nolimits_{n \geq 0} R^pf_*(\mathcal{I}^n\mathcal{F})
$$
is a graded $\mathcal{A} = \bigoplus_{n \geq 0} \mathcal{I}^n$-module
which is quasi-coherent and of finite type.
\end{lemma}

\begin{proof}
The statement is local on $Y$, hence this reduces to the
case where $Y$ is affine. In the affine case the result follows
from Lemma \ref{lemma-cohomology-powers-ideal-times-F}.
Details omitted.
\end{proof}

\begin{lemma}
\label{lemma-cohomology-powers-ideal-application}
Let $A$ be a Noetherian ring.
Let $I \subset A$ be an ideal.
Let $f : X \to \Spec(A)$ be a proper morphism.
Let $\mathcal{F}$ be a coherent sheaf on $X$.
Then for every $p \geq 0$ there exists an integer $c \geq 0$
such that
\begin{enumerate}
\item the multiplication map
$I^{n - c} \otimes H^p(X, I^c\mathcal{F}) \to H^p(X, I^n\mathcal{F})$
is surjective for all $n \geq c$,
\item the image of $H^p(X, I^{n + m}\mathcal{F}) \to H^p(X, I^n\mathcal{F})$
is contained in the submodule $I^{m - e} H^p(X, I^n\mathcal{F})$
where $e = \max(0, c - n)$ for $n + m \geq c$, $n, m \geq 0$,
\item we have
$$
\Ker(H^p(X, I^n\mathcal{F}) \to H^p(X, \mathcal{F})) =
\Ker(H^p(X, I^n\mathcal{F}) \to H^p(X, I^{n - c}\mathcal{F}))
$$
for $n \geq c$,
\item there are maps $I^nH^p(X, \mathcal{F}) \to H^p(X, I^{n - c}\mathcal{F})$
for $n \geq c$ such that the compositions
$$
H^p(X, I^n\mathcal{F}) \to 
I^{n - c}H^p(X, \mathcal{F}) \to
H^p(X, I^{n - 2c}\mathcal{F})
$$
and
$$
I^nH^p(X, \mathcal{F}) \to
H^p(X, I^{n - c}\mathcal{F}) \to
I^{n - 2c}H^p(X, \mathcal{F})
$$
for $n \geq 2c$ are the canonical ones, and
\item the inverse systems $(H^p(X, I^n\mathcal{F}))$ and
$(I^nH^p(X, \mathcal{F}))$ are pro-isomorphic.
\end{enumerate}
\end{lemma}

\begin{proof}
Write $M_n = H^p(X, I^n\mathcal{F})$ for $n \geq 1$ and
$M_0 = H^p(X, \mathcal{F})$ so that we have maps
$\ldots \to M_3 \to M_2 \to M_1 \to M_0$. Setting
$B = \bigoplus_{n \geq 0} I^n$, then $M = \bigoplus_{n \geq 0} M_n$
is a finite graded $B$-module, see
Lemma \ref{lemma-cohomology-powers-ideal-times-F}.
Observe that the products
$B_n \otimes M_m \to M_{m + n}$, $a \otimes m \mapsto a \cdot m$
are compatible with the maps in our inverse system in the sense
that the diagrams
$$
\xymatrix{
B_n \otimes_A M_m \ar[r] \ar[d] & M_{n + m} \ar[d] \\
B_n \otimes_A M_{m'} \ar[r] & M_{n + m'}
}
$$
commute for $n, m' \geq 0$ and $m \geq m'$.

\medskip\noindent
Proof of (1). Choose $d_1, \ldots, d_t \geq 0$ and $x_i \in M_{d_i}$
such that $M$ is generated by $x_1, \ldots, x_t$ over $B$.
For any $c \geq \max\{d_i\}$ we conclude that
$B_{n - c} \cdot M_c = M_n$ for $n \geq c$ and we conclude (1) is true.

\medskip\noindent
Proof of (2). Let $c$ be as in the proof of (1). Let $n + m \geq c$.
We have $M_{n + m} = B_{n + m - c} \cdot M_c$.
If $c > n$ then we use $M_c \to M_n$ and the compatibility of products
with transition maps pointed out above to conclude that
the image of $M_{n + m} \to M_n$ is contained in $I^{n + m - c}M_n$.
If $c \leq n$, then we write $M_{n + m} = B_m \cdot B_{n - c} \cdot M_c =
B_m \cdot M_n$ to see that the image is contained in $I^m M_n$.
This proves (2).

\medskip\noindent
Let $K_n \subset M_n$ be the kernel of the map $M_n \to M_0$. The
compatibility of products with transition maps pointed out above
shows that $K = \bigoplus K_n \subset M$ is a graded $B$-submodule.
As $B$ is Noetherian and $M$ is a finitely generated graded $B$-module,
this shows that $K$ is a finitely generated graded $B$-module.
Choose $d'_1, \ldots, d'_{t'} \geq 0$ and $y_i \in K_{d'_i}$
such that $K$ is generated by $y_1, \ldots, y_{t'}$ over $B$.
Set $c = \max(d'_i, d'_j)$. Since $y_i \in \Ker(M_{d'_i} \to M_0)$
we see that $B_n \cdot y_i \subset \Ker(M_{n + d'_i} \to M_n)$.
In this way we see that $K_n = \Ker(M_n \to M_{n - c})$ for $n \geq c$.
This proves (3).

\medskip\noindent
Consider the following commutative solid diagram
$$
\xymatrix{
I^n \otimes_A M_0 \ar[r] \ar[d] &
I^nM_0 \ar[r] \ar@{..>}[d] &
M_0 \ar[d] \\
M_n \ar[r] &
M_{n - c} \ar[r] &
M_0
}
$$
Since the kernel of the surjective arrow $I^n \otimes_A M_0 \to I^nM_0$
maps into $K_n$ by the above we obtain the dotted arrow and the
composition $I^nM_0 \to M_{n - c} \to M_0$ is the canonical map.
Then clearly the composition $I^nM_0 \to M_{n - c} \to I^{n - 2c}M_0$
is the canonical map for $n \geq 2c$. Consider the composition
$M_n \to I^{n - c}M_0 \to M_{n - 2c}$. The first map
sends an element of the form $a \cdot m$
with $a \in I^{n - c}$ and $m \in M_c$
to $a m'$ where $m'$ is the image of $m$ in $M_0$.
Then the second map sends this to $a \cdot m'$ in $M_{n - 2c}$ and
we see (4) is true.

\medskip\noindent
Part (5) is an immediate consequence of (4) and the definition of
morphisms of pro-objects.
\end{proof}

\noindent
In the situation of Lemmas \ref{lemma-cohomology-powers-ideal-times-F} and
\ref{lemma-cohomology-powers-ideal-application} consider the inverse
system
$$
\mathcal{F}/I\mathcal{F} \leftarrow
\mathcal{F}/I^2\mathcal{F} \leftarrow
\mathcal{F}/I^3\mathcal{F} \leftarrow \ldots
$$
We would like to know what happens to the cohomology groups.
Here is a first result.

\begin{lemma}
\label{lemma-ML-cohomology-powers-ideal}
Let $A$ be a Noetherian ring.
Let $I \subset A$ be an ideal.
Let $f : X \to \Spec(A)$ be a proper morphism.
Let $\mathcal{F}$ be a coherent sheaf on $X$.
Fix $p \geq 0$. There exists a $c \geq 0$ such that
\begin{enumerate}
\item for all $n \geq c$ we have
$$
\Ker(H^p(X, \mathcal{F}) \to H^p(X, \mathcal{F}/I^n\mathcal{F})) \subset
I^{n - c}H^p(X, \mathcal{F}).
$$
\item the inverse system
$$
\left(H^p(X, \mathcal{F}/I^n\mathcal{F})\right)_{n \in \mathbf{N}}
$$
satisfies the Mittag-Leffler condition (see
Homology, Definition \ref{homology-definition-Mittag-Leffler}), and
\item we have
$$
\Im(H^p(X, \mathcal{F}/I^k\mathcal{F})
\to H^p(X, \mathcal{F}/I^n\mathcal{F}))
=
\Im(H^p(X, \mathcal{F})
\to H^p(X, \mathcal{F}/I^n\mathcal{F}))
$$
for all $k \geq n + c$.
\end{enumerate}
\end{lemma}

\begin{proof}
Let $c = \max\{c_p, c_{p + 1}\}$, where $c_p, c_{p + 1}$ are the integers
found in Lemma \ref{lemma-cohomology-powers-ideal-application} for
$H^p$ and $H^{p + 1}$.

\medskip\noindent
Let us prove part (1). Consider the short exact sequence
$$
0 \to I^n\mathcal{F} \to \mathcal{F} \to \mathcal{F}/I^n\mathcal{F} \to 0
$$
From the long exact cohomology sequence we see that
$$
\Ker(
H^p(X, \mathcal{F}) \to H^p(X, \mathcal{F}/I^n\mathcal{F})
)
=
\Im(
H^p(X, I^n\mathcal{F}) \to H^p(X, \mathcal{F})
)
$$
Hence by Lemma \ref{lemma-cohomology-powers-ideal-application} part (2)
we see that this is contained in $I^{n - c}H^p(X, \mathcal{F})$ for $n \geq c$.

\medskip\noindent
Note that part (3) implies part (2) by definition of the Mittag-Leffler
systems.

\medskip\noindent
Let us prove part (3). Fix an $n$. Consider the commutative diagram
$$
\xymatrix{
0 \ar[r] &
I^n\mathcal{F} \ar[r] &
\mathcal{F} \ar[r] &
\mathcal{F}/I^n\mathcal{F} \ar[r] &
0 \\
0 \ar[r] &
I^{n + m}\mathcal{F} \ar[r] \ar[u] &
\mathcal{F} \ar[r] \ar[u] &
\mathcal{F}/I^{n + m}\mathcal{F} \ar[r] \ar[u] &
0
}
$$
This gives rise to the following commutative diagram
$$
\xymatrix{
H^p(X, \mathcal{F}) \ar[r] &
H^p(X, \mathcal{F}/I^n\mathcal{F}) \ar[r]_\delta &
H^{p + 1}(X, I^n\mathcal{F}) \ar[r] &
H^{p + 1}(X, \mathcal{F}) \\
H^p(X, \mathcal{F}) \ar[r] \ar[u]^1 &
H^p(X, \mathcal{F}/I^{n + m}\mathcal{F}) \ar[r] \ar[u]^\gamma &
H^{p + 1}(X, I^{n + m}\mathcal{F}) \ar[u]^\alpha \ar[r]^-\beta &
H^{p + 1}(X, \mathcal{F}) \ar[u]_1
}
$$
with exact rows. By
Lemma \ref{lemma-cohomology-powers-ideal-application} part (4) the kernel
of $\beta$ is equal to the kernel of $\alpha$ for $m \geq c$.
By a diagram chase this shows that the image of $\gamma$ is contained
in the kernel of $\delta$ which shows that part (3) is true
(set $k = n + m$ to get it).
\end{proof}

\begin{theorem}[Theorem on formal functions]
\label{theorem-formal-functions}
Let $A$ be a Noetherian ring.
Let $I \subset A$ be an ideal.
Let $f : X \to \Spec(A)$ be a proper morphism.
Let $\mathcal{F}$ be a coherent sheaf on $X$.
Fix $p \geq 0$.
The system of maps
$$
H^p(X, \mathcal{F})/I^nH^p(X, \mathcal{F})
\longrightarrow
H^p(X, \mathcal{F}/I^n\mathcal{F})
$$
define an isomorphism of limits
$$
H^p(X, \mathcal{F})^\wedge
\longrightarrow
\lim_n H^p(X, \mathcal{F}/I^n\mathcal{F})
$$
where the left hand side is the completion of the $A$-module
$H^p(X, \mathcal{F})$ with respect to the ideal $I$, see
Algebra, Section \ref{algebra-section-completion}.
Moreover, this is in fact a homeomorphism for the limit topologies.
\end{theorem}

\begin{proof}
This follows from Lemma \ref{lemma-ML-cohomology-powers-ideal} as follows.
Set $M = H^p(X, \mathcal{F})$, $M_n = H^p(X, \mathcal{F}/I^n\mathcal{F})$,
and denote $N_n = \Im(M \to M_n)$. By
Lemma \ref{lemma-ML-cohomology-powers-ideal} parts (2) and (3) we see that
$(M_n)$ is a Mittag-Leffler system with
$N_n \subset M_n$ equal to the image of $M_k$ for all $k \gg n$.
It follows that $\lim M_n = \lim N_n$ as topological modules (with limit
topologies). On the other hand, the $N_n$ form an inverse system of quotients
of the module $M$ and hence $\lim N_n$ is the completion of $M$
with respect to the topology given by the kernels $K_n = \Ker(M \to N_n)$.
By Lemma \ref{lemma-ML-cohomology-powers-ideal} part (1)
we have $K_n \subset I^{n - c}M$ and since $N_n \subset M_n$
is annihilated by $I^n$ we have $I^n M \subset K_n$.
Thus the topology defined using the submodules $K_n$
as a fundamental system of open neighbourhoods of $0$
is the same as the $I$-adic topology and we find that the
induced map $M^\wedge = \lim M/I^nM \to \lim N_n = \lim M_n$
is an isomorphism of topological modules\footnote{
To be sure, the limit topology on $M^\wedge$ is the same as its
$I$-adic topology as follows from
Algebra, Lemma \ref{algebra-lemma-hathat-finitely-generated}.
See More on Algebra, Section \ref{more-algebra-section-topological-ring}.}.
\end{proof}

\begin{lemma}
\label{lemma-spell-out-theorem-formal-functions}
Let $A$ be a ring. Let $I \subset A$ be an ideal. Assume $A$ is
Noetherian and complete with respect to $I$.
Let $f : X \to \Spec(A)$ be a proper morphism.
Let $\mathcal{F}$ be a coherent sheaf on $X$.
Then
$$
H^p(X, \mathcal{F}) = \lim_n H^p(X, \mathcal{F}/I^n\mathcal{F})
$$
for all $p \geq 0$.
\end{lemma}

\begin{proof}
This is a reformulation of the theorem on formal functions
(Theorem \ref{theorem-formal-functions}) in the
case of a complete Noetherian base ring. Namely, in this case the
$A$-module $H^p(X, \mathcal{F})$ is finite
(Lemma \ref{lemma-proper-over-affine-cohomology-finite}) hence
$I$-adically complete (Algebra, Lemma \ref{algebra-lemma-completion-tensor})
and we see that completion on the left hand side is not necessary.
\end{proof}

\begin{lemma}
\label{lemma-formal-functions-stalk}
Given a morphism of schemes $f : X \to Y$ and a quasi-coherent sheaf
$\mathcal{F}$ on $X$. Assume
\begin{enumerate}
\item $Y$ locally Noetherian,
\item $f$ proper, and
\item $\mathcal{F}$ coherent.
\end{enumerate}
Let $y \in Y$ be a point. Consider the infinitesimal neighbourhoods
$$
\xymatrix{
X_n =
\Spec(\mathcal{O}_{Y, y}/\mathfrak m_y^n) \times_Y X
\ar[r]_-{i_n} \ar[d]_{f_n} &
X \ar[d]^f \\
\Spec(\mathcal{O}_{Y, y}/\mathfrak m_y^n) \ar[r]^-{c_n} & Y
}
$$
of the fibre $X_1 = X_y$ and set $\mathcal{F}_n = i_n^*\mathcal{F}$.
Then we have
$$
\left(R^pf_*\mathcal{F}\right)_y^\wedge
\cong
\lim_n H^p(X_n, \mathcal{F}_n)
$$
as $\mathcal{O}_{Y, y}^\wedge$-modules.
\end{lemma}

\begin{proof}
This is just a reformulation of a special case of the theorem
on formal functions, Theorem \ref{theorem-formal-functions}.
Let us spell it out. Note that $\mathcal{O}_{Y, y}$ is a Noetherian
local ring. Consider the canonical morphism
$c : \Spec(\mathcal{O}_{Y, y}) \to Y$, see
Schemes, Equation (\ref{schemes-equation-canonical-morphism}).
This is a flat morphism as it identifies local rings.
Denote momentarily $f' : X' \to \Spec(\mathcal{O}_{Y, y})$
the base change of $f$ to this local ring. We see that
$c^*R^pf_*\mathcal{F} = R^pf'_*\mathcal{F}'$ by
Lemma \ref{lemma-flat-base-change-cohomology}.
Moreover, the infinitesimal neighbourhoods of
the fibre $X_y$ and $X'_y$ are identified (verification omitted; hint:
the morphisms $c_n$ factor through $c$).

\medskip\noindent
Hence we may assume that $Y = \Spec(A)$ is the spectrum of
a Noetherian local ring $A$ with maximal ideal $\mathfrak m$
and that $y \in Y$ corresponds to the closed point (i.e., to $\mathfrak m$).
In particular it follows that
$$
\left(R^pf_*\mathcal{F}\right)_y =
\Gamma(Y, R^pf_*\mathcal{F}) =
H^p(X, \mathcal{F}).
$$

\medskip\noindent
In this case also, the morphisms $c_n$ are each closed immersions.
Hence their base changes $i_n$ are closed immersions as well.
Note that $i_{n, *}\mathcal{F}_n = i_{n, *}i_n^*\mathcal{F}
= \mathcal{F}/\mathfrak m^n\mathcal{F}$. By the Leray spectral sequence
for $i_n$, and Lemma \ref{lemma-finite-pushforward-coherent} we see that
$$
H^p(X_n, \mathcal{F}_n) =
H^p(X, i_{n, *}\mathcal{F}_n) =
H^p(X, \mathcal{F}/\mathfrak m^n\mathcal{F})
$$
Hence we may indeed apply the theorem on formal functions to compute
the limit in the statement of the lemma and we win.
\end{proof}

\noindent
Here is a lemma which we will generalize later to fibres of
dimension $ > 0$, namely the next lemma.

\begin{lemma}
\label{lemma-higher-direct-images-zero-finite-fibre}
Let $f : X \to Y$ be a morphism of schemes.
Let $y \in Y$.
Assume
\begin{enumerate}
\item $Y$ locally Noetherian,
\item $f$ is proper, and
\item $f^{-1}(\{y\})$ is finite.
\end{enumerate}
Then for any coherent sheaf $\mathcal{F}$ on $X$ we have
$(R^pf_*\mathcal{F})_y = 0$ for all $p > 0$.
\end{lemma}

\begin{proof}
The fibre $X_y$ is finite, and by
Morphisms, Lemma \ref{morphisms-lemma-finite-fibre} it
is a finite discrete space. Moreover, the underlying topological
space of each infinitesimal neighbourhood $X_n$ is the same.
Hence each of the schemes $X_n$ is affine according to
Schemes, Lemma \ref{schemes-lemma-scheme-finite-discrete-affine}.
Hence it follows that $H^p(X_n, \mathcal{F}_n) = 0$ for all
$p > 0$. Hence we see that $(R^pf_*\mathcal{F})_y^\wedge = 0$
by Lemma \ref{lemma-formal-functions-stalk}.
Note that $R^pf_*\mathcal{F}$ is coherent by
Proposition \ref{proposition-proper-pushforward-coherent} and
hence $R^pf_*\mathcal{F}_y$ is a finite
$\mathcal{O}_{Y, y}$-module. By Nakayama's lemma
(Algebra, Lemma \ref{algebra-lemma-NAK})
if the completion of a finite module over a local ring
is zero, then the module is zero. Whence
$(R^pf_*\mathcal{F})_y = 0$.
\end{proof}

\begin{lemma}
\label{lemma-higher-direct-images-zero-above-dimension-fibre}
Let $f : X \to Y$ be a morphism of schemes.
Let $y \in Y$.
Assume
\begin{enumerate}
\item $Y$ locally Noetherian,
\item $f$ is proper, and
\item $\dim(X_y) = d$.
\end{enumerate}
Then for any coherent sheaf $\mathcal{F}$ on $X$ we have
$(R^pf_*\mathcal{F})_y = 0$ for all $p > d$.
\end{lemma}

\begin{proof}
The fibre $X_y$ is of finite type over $\Spec(\kappa(y))$.
Hence $X_y$ is a Noetherian scheme by
Morphisms, Lemma \ref{morphisms-lemma-finite-type-noetherian}.
Hence the underlying topological space of $X_y$ is Noetherian, see
Properties, Lemma \ref{properties-lemma-Noetherian-topology}.
Moreover, the underlying topological space of each infinitesimal
neighbourhood $X_n$ is the same as that of $X_y$.
Hence $H^p(X_n, \mathcal{F}_n) = 0$ for all $p > d$ by
Cohomology, Proposition \ref{cohomology-proposition-vanishing-Noetherian}.
Hence we see that $(R^pf_*\mathcal{F})_y^\wedge = 0$
by Lemma \ref{lemma-formal-functions-stalk} for $p > d$.
Note that $R^pf_*\mathcal{F}$ is coherent by
Proposition \ref{proposition-proper-pushforward-coherent} and
hence $R^pf_*\mathcal{F}_y$ is a finite
$\mathcal{O}_{Y, y}$-module. By Nakayama's lemma
(Algebra, Lemma \ref{algebra-lemma-NAK})
if the completion of a finite module over a local ring
is zero, then the module is zero. Whence
$(R^pf_*\mathcal{F})_y = 0$.
\end{proof}








